\section{Introduction}\label{sec:introduction}
At the critical temperature of the Sherrington--Kirkpatrick model,
equilibrium fluctuations contain a direction that local spin updates
can change only slowly. We use this fluctuation to prove a total-variation
obstruction for most equilibrium configurations viewed as fixed initial
states. The same static estimates determine the covariance scale and
the relaxation detected by linear observables.

\subsection{The model and main results}
The model introduced by Sherrington and Kirkpatrick \cite{sk1975}
places spins on $\Omega_n=\{-1,1\}^n$ and couples every pair through an
independent Gaussian variable. At inverse temperature $\beta=1$ and zero field,
$W_{ij}\sim N(0,1/n)$ for $i<j$, and
\begin{equation}\label{eq:intro-law}
 \mu_W(x)\propto
 \exp\!\left(\sum_{i<j}W_{ij}x_ix_j\right),
 \qquad x\in\{-1,1\}^n.
\end{equation}
Here $\beta=1$ marks the large-$n$ equilibrium spin-glass transition
\cite[Section~1, after Equation~(1.4)]{duHuang2026FreeEnergy}. The overlap
$n^{-1}x^Ty$ of two independent equilibrium replicas has scale $n^{-1/3}$
at criticality, compared with its Gaussian $n^{-1/2}$ scale at each fixed
$\beta<1$ \cite[Equation~(1.1) and Theorem~1.1]{duHuang2026Overlap}.
A heat-bath update resamples one spin from its conditional law under
$\mu_W$. We use two clocks: in continuous time each site has an independent
rate-one clock; in discrete time each step attempts one update at a
uniformly chosen site. An attempt may leave the spin unchanged.

For fixed disorder and initial configuration $x$, let
$d_t^{\rm ct}(W,x)$ and $d_k^{\rm dt}(W,x)$ denote the total-variation
distances of the corresponding transition laws from $\mu_W$.
Section~\ref{sec:model} gives the operators and conventions in detail.
Our first result concerns an equilibrium configuration after its realized
value has been fixed as the initial state. All probability limits in the
main results are over the original off-diagonal Gaussian disorder.

\begin{theorem}[Realized equilibrium initial states]\label{thm:gibbs-start}
For every deterministic sequence $t_n\ge0$ with $t_n=o(n^{2/3})$,
\[
 \mu_W\{x:d_{t_n}^{\rm ct}(W,x)>1/4\}
 \xrightarrow{\Prob}1.
\]
For every deterministic integer sequence $k_n\ge0$ with $k_n=o(n^{5/3})$,
\[
 \mu_W\{x:d_{k_n}^{\rm dt}(W,x)>1/4\}
 \xrightarrow{\Prob}1.
\]
The convergence is in probability over the disorder.
\end{theorem}

The theorem tests each realized start separately: averaging its transition
law over the equilibrium initial configuration would give the stationary
law $\mu_W$. It also treats each prescribed deterministic time sequence
separately; the set of slowly mixing starts may depend on that sequence.

To state the static conclusion, set $\Sigma=\Cov_{\mu_W}(x)$ and define
the unscaled conditional-variance form
\[
 \mathcal D_{\mu_W}(f)
 =\sum_{i=1}^n\mu_W[\Var(f\mid x_{-i})].
\]
Here $x_{-i}$ comprises every spin except $i$. This is the Dirichlet
form for the rate-one-per-site clock. The largest inverse Rayleigh quotient
among linear functions is
\[
 \mathcal R_{\rm lin}(\mu_W)
 =\sup_{a\in\R^n\setminus\{0\}}
   \frac{\Var_{\mu_W}(a^Tx)}{\mathcal D_{\mu_W}(a^Tx)}.
\]

\begin{theorem}[Covariance and linear tests]\label{thm:linear-scale}
For every positive deterministic sequence $M_n\to\infty$, with probability tending to one both $\|\Sigma\|$ and $\mathcal R_{\rm lin}(\mu_W)$ belong to
\[
 \left[\frac{n^{2/3}}{M_n},\ M_n n^{2/3}\right].
\]
\end{theorem}

The factor $M_n$ may diverge more slowly than any prescribed power or
logarithm. Because the quotient here ranges only over linear functions,
its upper bound does not control the full inverse spectral gap.

For comparison with customary worst-start mixing, let $t_{\rm mix}$ and
$k_{\rm mix}$ be the first continuous time and the first integer step,
respectively, at which every initial configuration has distance at most
$1/4$. Theorem~\ref{thm:gibbs-start} gives the lower bounds below. An
independent comparison with one fixed high-temperature model gives the
subexponential upper bounds.

\begin{theorem}[Critical mixing bounds]\label{thm:mixing}
For every fixed $\epsilon>0$,
\begin{align}
 \Prob\!\left\{n^{2/3-\epsilon}\le t_{\rm mix}
                   \le e^{\epsilon n}\right\}&\longrightarrow1,
       \label{eq:continuous-main}\\
 \Prob\!\left\{n^{5/3-\epsilon}\le k_{\rm mix}
                   \le e^{\epsilon n}\right\}&\longrightarrow1.
       \label{eq:discrete-main}
\end{align}
\end{theorem}

\subsection{Historical and mathematical context}
Critical spin-glass relaxation has been studied through several models
and convergence criteria. Kirkpatrick and Sherrington
\cite{kirkpatrickSherrington1978} investigated Glauber autocorrelations
above the critical temperature using a linearized mean-field treatment.
Sompolinsky and Zippelius developed equations for disorder-averaged
correlation and response \cite{sompolinskyZippelius1981}, and treated
soft-spin Langevin dynamics in \cite{sompolinskyZippelius1982}. For Ising dynamics,
Billoire and Campbell \cite{billoireCampbell2011} numerically examined
equilibrium autocorrelations on an $n^{2/3}$ scale, with time in updates
per spin. These studies concern correlation and response; the present
paper instead tests total variation after each realized equilibrium start.

Our static comparison uses a spherical model, a classical counterpart of
the Ising spin glass. In their thermodynamic analysis, Kosterlitz,
Thouless, and Jones \cite{kosterlitzThoulessJones1976} located its transition
at the spectral edge and identified the top eigenmode as the ordering
direction. Comets
\cite[Equations~(2.1)--(2.2)]{comets1996} gave the fixed-spectrum
Haar-averaging identity relating the two partition functions. Baik and Lee
\cite{baikLee2016} proved Gaussian and Tracy--Widom spherical free-energy
limits at fixed temperatures on the two sides of criticality, and Landon
\cite{landon2022} treated the critical Gaussian limit and the saddle's
$n^{-2/3}$ edge scale. Bhattacharya and Sen \cite{bhattacharyaSen2016}
reduced fixed-spectrum Haar second moments through the two eigenvalues
$1\pm q$, where $q=x^Ty/n$ is the overlap of two vectors of norm $\sqrt n$,
comparing free-energy exponential rates under
high-temperature spectral assumptions.

The critical results of Du and Huang are close predecessors of the static
estimates used here. For the normalized Ising and spherical partition
functions $Z_{\mathrm{Is}}$ and $Z_{\mathrm{sph}}$ built from the same GOE
matrix, their free-energy work
\cite[Theorems~1.4 and~1.6]{duHuang2026FreeEnergy} proves
\[
 \E\!\left[\left(\frac{Z_{\mathrm{Is}}}{Z_{\mathrm{sph}}}-1\right)^2\right]
 =O(n^{-1/3})
\]
and an annealed squared-overlap estimate of order $n^{-2/3}$.
Their subsequent work \cite[Theorem~1.1]{duHuang2026Overlap}
determines the limiting quenched critical overlap distribution.

Those results also give an alternative route to the static inputs of
this paper. Proposition~4.1 in \cite{duHuang2026Overlap} gives convergence
of the rescaled spectral coefficients, while the proof of Lemma~4.5
identifies the corresponding spherical coordinates. Together with the
fixed-coordinate convergence in the proof of Lemma~4.7,
Equations~(4.24)--(4.26), and the absolutely continuous limiting marginal
in Equation~(2.14), these facts imply
vanishing mass in every deterministic shrinking top-coordinate window.
The restricted Haar first moment and the partition comparison in
\cite[Theorem~1.6]{duHuang2026FreeEnergy} then transfer this estimate to
the cube. The annealed squared-overlap bound also implies, by Markov's inequality, that
$n^{2/3}\E_{\mu_W^{\otimes2}}q(x,y)^2$ is bounded in probability; this is
the overlap input for the covariance upper bound and the cavity argument.
The sign and cavity arguments in Section~\ref{sec:dynamics} give the
additional dynamical and linear-observable deductions.
Our technical sections provide a separate static proof, uniform on
chosen deterministic spectral events, using exact Gaussian conditioning,
a quantitative cubic deficit, and monotone outward overlap bins.

Two further groups of methods enter the proof. The GOE tridiagonal
representation is classical \cite{dumitriuEdelman2002}; our spectral
estimates develop the quadratic-form approach of Ram\'irez, Rider, and
Vir\'ag \cite{ramirezRiderVirag2011} and Ledoux and Rider
\cite{ledouxRider2010}. For the independent upper comparison we use the
fixed-$\beta<1$, zero-field heat-bath Poincar\'e theorem
\cite[Theorem~1.1]{skgapcompanion}. Bauerschmidt and Bodineau
\cite{bauerschmidtBodineau2019} proved a logarithmic Sobolev inequality
with the full spin-flip-gradient form, while Eldan, Koehler, and Zeitouni
\cite{eldanKoehlerZeitouni2022} proved a heat-bath Poincar\'e inequality
for interaction matrices satisfying $0\preceq J\prec I$. After a diagonal shift, their SK applications
reach $\beta<1/4$. More recently, Wang \cite{wang2026SKMixing} proved
worst-start heat-bath mixing to total-variation error $\varepsilon$ in
$O_\beta(n\log(n/\varepsilon))$ attempted updates for fixed $\beta<1/2$,
uniformly over external fields on one
high-probability disorder event. Boban, Li, and Oveis Gharan
\cite{bobanLiOveisGharan2026SK} proved a constant unscaled gap and
zero-field worst-start mixing in $O(n^2)$ attempted updates for $\beta<1/2+\varepsilon_0$
with an absolute $\varepsilon_0>0$, with high probability over the disorder.
The fixed-$\beta<1$ companion input
used here covers the temperatures required by the comparison in
Section~\ref{sec:dynamics}, which fixes $\beta$ before taking $n$ to infinity.

The companion \emph{Critical mixing in the Sherrington--Kirkpatrick model}
\cite[Theorem~1.1]{criticalmixingcompanion} proves matching mixing exponents
in probability over the disorder: $n^{2/3+o(1)}$ in continuous time and
$n^{5/3+o(1)}$ attempted updates, together with exponent $2/3$ for the
full relaxation and inverse classical logarithmic Sobolev constants.
Each exponent statement allows every fixed positive power slack. Its
lower and static inputs are the realized-start and covariance/linear
results in Theorems~\ref{thm:gibbs-start} and~\ref{thm:linear-scale} here;
its upper estimates use separate arguments. The independent comparison
in Theorem~\ref{thm:mixing} remains part of the present paper.

\subsection{How the proof works}
Following the augmentation used by Comets \cite[Equation~(1.1)]{comets1996},
add an independent Gaussian diagonal to the coupling matrix. Its energy
is constant on the cube, so the Gibbs law and both chains stay the same,
while the augmented matrix is GOE. Let $v$ be a unit top eigenvector of
that matrix. The main static input is
\[
 \mu_W\{|v^Tx|\le b_n n^{1/3}\}\longrightarrow0
 \quad\text{in probability for every positive deterministic }b_n\to0.
\]
Thus the equilibrium projection is rarely small on the $n^{1/3}$ scale.
In contrast, $\mathcal D_{\mu_W}(v^Tx)\le1$. Reversibility bounds the
mean-square change along a stationary trajectory by $2t$ in continuous
time and $2k/n$ after $k$ attempted updates. At each prescribed
deterministic time that is little-oh of the corresponding $n^{2/3}$ or
$n^{5/3}$ scale, the two endpoint projections have the same nonzero sign
with high probability. For each fixed initial state, that sign defines a
half-space whose equilibrium mass is at most one half, giving the
total-variation test in Theorem~\ref{thm:gibbs-start}.

The scale of the projection comes from the spectral edge. Three static
stages make this connection precise.
Section~\ref{sec:spectral} controls the leading eigenvalues and two
resolvent traces at distance $s_0=L_n n^{-2/3}$ above the GOE edge, on
an event depending only on eigenvalues. Here $L_n$ may diverge arbitrarily
slowly. Section~\ref{sec:spherical} writes the tilted spherical law as
follows: for the ordered eigenvalues $\lambda_i$, set
$\Lambda=\diag(\lambda_1,\ldots,\lambda_n)$ and tilt uniform measure on
$\{u\in\R^n:\|u\|^2=n\}$ by $\exp(u^T\Lambda u/2)$. This law is represented by
independent Gaussian coordinates conditioned on their squared norm.
The edge estimate gives the leading coordinate a Gaussian variance at
least a constant times $s_0^{-1}$, so its standard deviation is at least
a constant times $n^{1/3}/\sqrt{L_n}$. Lower and upper radial-density
bounds control the effect of conditioning, giving small-projection and
squared-overlap estimates on the sphere.

Section~\ref{sec:orientation} transfers these estimates to the cube.
Haar averaging identifies the first moment of an Ising partition sum with
the spherical integral. To control normalized Gibbs probabilities, the
section also proves concentration of that partition sum. Its second
moment integrates the same pair integral against two overlap laws: a
binomial law on the cube and a continuous law on the sphere. The spectral
resolvent and spherical saddle estimates give a cubic deficit that
suppresses large overlaps. For small overlaps, disjoint outward bins and
monotonicity compare the discrete sum with the spherical integral without
differentiating the pair integral. The first-moment identity and
partition concentration then transfer small projections, while the
weighted second moment transfers the squared-overlap bound.

These two outputs give opposite sides of the covariance estimate: the
top projection gives a lower bound, and the squared overlap controls
$\Tr\Sigma^2$ and hence gives an upper bound. For linear
inverse Rayleigh quotients, covariance alone is not enough, since the
Dirichlet denominator weights each coefficient by a mean one-site conditional
variance. Section~\ref{sec:dynamics} uses spin deletion to show that sites
with large reciprocal weights contribute only a smaller covariance block.
The same section proves the sign test, the independent high-temperature
upper comparison, and the three main results. Choosing $L_n$ after each
prescribed deterministic threshold preserves the quantifiers in the
realized-start and static results.
Figure~\ref{fig:proof-structure} summarizes the lower and static argument;
Appendix~\ref{sec:specification} gives elementary examples clarifying the
role of the update rule, clock, and disorder quantifiers.

\begin{figure}[t]
\centering
\begin{tikzpicture}[
 box/.style={draw=blue!35!black,rounded corners=2pt,fill=blue!3,
  text width=7.2cm,align=center,inner sep=5pt,font=\small},
 branch/.style={box,text width=5cm},
 flow/.style={-{Latex[length=2mm]},semithick,draw=blue!50!black}]
\node[box] (spec) {GOE edge and resolvent control\\Proposition~\ref{prop:spectral}};
\node[box,below=8mm of spec] (sphere) {Spherical estimates from Gaussian conditioning\\Partition, projection, and overlap bounds\\Section~\ref{sec:spherical}};
\node[box,below=8mm of sphere] (cube) {Haar first and second moments\\Partition concentration and observable transfer\\Section~\ref{sec:orientation}};
\node[branch,anchor=north] (small) at ([xshift=-2.8cm,yshift=-12mm]cube.south) {Small-projection mass on the cube\\Theorem~\ref{thm:static-small-ball}};
\node[branch,anchor=north] (cov) at ([xshift=2.8cm,yshift=-12mm]cube.south) {Squared-overlap control on the cube\\Proposition~\ref{prop:overlap}};
\node[branch,below=8mm of small] (slow) {Sign test for realized initial states\\Theorem~\ref{thm:gibbs-start}};
\node[branch,below=8mm of cov] (lin) {Covariance scale and linear\\inverse Rayleigh bounds\\Theorem~\ref{thm:linear-scale}};
\draw[flow] (spec)--(sphere);
\draw[flow] (sphere)--(cube);
\draw[flow] (spec.west)--++(-6mm,0)|-(cube.west);
\draw[flow] (cube.south)--++(0,-.6)-|(small.north);
\draw[flow] (cube.south)--++(0,-.6)-|(cov.north);
\draw[flow] (small)--(slow);
\draw[flow] (cov)--(lin);
\draw[flow] (small.east)--(lin.west);
\end{tikzpicture}
\caption{The lower and static argument. The direct arrow from the spectral estimates to Haar comparison records their additional use in the cubic overlap deficit. The restricted first moment transfers the small-projection numerator, while partition concentration controls its normalization; a weighted second moment transfers squared overlap. Small projections give the covariance lower bound and squared overlap gives its upper bound. A separate cavity estimate then controls the linear inverse Rayleigh quotients. The independent subexponential upper comparison is outside the diagram.}
\label{fig:proof-structure}
\end{figure}

